\documentclass[10pt,a4paper]{article}

\usepackage[utf8]{inputenc}
\usepackage[T1]{fontenc}
\usepackage{lmodern}
\usepackage[italian]{babel}
\usepackage[a4paper,margin=2.05cm]{geometry}
\usepackage{amsmath,amssymb}
\usepackage{booktabs}
\usepackage{graphicx}
\usepackage{abstract}
\renewcommand{\abstractnamefont}{\Large\bfseries}
\usepackage{float}
\usepackage{tikz}
\usepackage[ruled,vlined,linesnumbered]{algorithm2e}
\usepackage{caption}
\usepackage{microtype}
\usepackage[hidelinks]{hyperref}
\usetikzlibrary{arrows.meta,positioning,shapes.geometric,calc}
% Definizione colori per la copertina
\definecolor{unictblue}{RGB}{11, 25, 46}      % Blu scuro istituzionale
\definecolor{accentcyan}{RGB}{0, 136, 145}    % Ciano per accenti tecnici

\usepackage[compact]{titlesec}
\titlespacing*{\section}{0pt}{7pt plus 2pt minus 2pt}{3pt plus 1pt minus 1pt}
\titlespacing*{\subsection}{0pt}{5pt plus 1pt minus 1pt}{2pt}
\titlespacing*{\paragraph}{0pt}{4pt plus 1pt minus 1pt}{6pt}

\captionsetup{font=small,labelfont=bf}
\setlength{\parskip}{1pt}
\setlength{\abovecaptionskip}{4pt}
\setlength{\belowcaptionskip}{0pt}
\setlength{\textfloatsep}{9pt plus 2pt minus 2pt}
\setlength{\intextsep}{8pt plus 2pt minus 2pt}
\setlength{\floatsep}{7pt plus 2pt minus 2pt}


\begin{document}
% =====================================================================
%  COPERTINA (TITLE PAGE)
% =====================================================================
\begin{titlepage}
    \centering
    % Intestazione Università e Dipartimento
    {\scshape\Large Università degli Studi di Catania} \\
    \vspace{0.2cm}
    {\scshape\large Dipartimento di Matematica e Informatica} \\
    \vspace{0.15cm}
    {\small Corso di Laurea Magistrale in Informatica LM-18} \\
    \vspace{0.4cm}
    
    \vspace{7.5cm}
    
    % Linea decorativa superiore
    \color{unictblue}\rule{\textwidth}{1.5pt} \\[0.5cm]
    
    % Titolo Principale
    {\huge \bfseries Capacitated Vehicle Routing Problem} \\[0.4cm]
    
    % Sottotitolo
    {\large \bfseries \color{accentcyan} Un algoritmo immunologico a selezione clonale\\ ibridato con ricerca locale} \\[0.2cm]
    
    % Linea decorativa inferiore
    \color{unictblue}\rule{\textwidth}{1.5pt} \\[1.5cm]
    
    % Tipologia di documento
    {\large \scshape Relazione di Progetto --- Heuristics \& Metaheuristics for Optimization \& Learning} \\
    
    \vfill

    \vspace{4.5cm}
    % Sezione Autore e Docenti organizzata su due colonne pulite
    \begin{minipage}[t]{0.45\textwidth}
        \begin{flushleft} \large
            \emph{Studente:} \\
            \textbf{Massimiliano Cassia} \\
            \small Matr. \textit{100016487}
        \end{flushleft}
    \end{minipage}
    \hfill
    \begin{minipage}[t]{0.45\textwidth}
        \begin{flushright} \large
            \emph{Docente:} \\
            \textbf{Prof. Mario Pavone} 
        \end{flushright}
    \end{minipage}
    
    \vfill
    \vspace{1.5cm}
    
    % Piè di pagina con l'Anno Accademico corrente
    \color{black}
    {\large Anno Accademico 2025/2026}
    
\end{titlepage}
\pagenumbering{roman}
\begin{abstract}
\large
\noindent
Questo lavoro presenta un algoritmo immunologico ispirato al principio della
selezione clonale (opzione~2 della consegna), ibridato con una ricerca locale a
vicinati multipli, per il \emph{Capacitated Vehicle Routing Problem}. L'elemento
portante è la rappresentazione di ogni soluzione come singola permutazione di
clienti (\emph{giant tour}), con la suddivisione in percorsi delegata a un
decodificatore \emph{esatto} basato su programmazione dinamica (\emph{Split}).
Sulle dieci istanze CVRPLIB richieste, con $5$~esecuzioni indipendenti ciascuna e
criterio di arresto $\mathrm{FE}=3.5\times 10^5$, l'algoritmo raggiunge uno scarto
medio dall'ottimo noto dell'$1{,}82\%$ (minimo $0{,}17\%$, massimo $5{,}58\%$);
tutte le città sono servite in tutte le esecuzioni e il vincolo sul numero di
veicoli è sempre rispettato. Un'analisi conclusiva scompone in modo esatto lo
scarto residuo, mostrando che esso è interamente imputabile all'\emph{assegnamento}
dei clienti ai veicoli e non all'\emph{ordinamento} interno dei percorsi, che
risulta ottimo in tutte le soluzioni prodotte.
\end{abstract}
\newpage
\pagenumbering{arabic}
\tableofcontents
\bigskip
\newpage
\section{Introduzione}
\label{sec:intro}

Il \emph{Vehicle Routing Problem} (VRP) è uno dei problemi
classici dell'ottimizzazione combinatoria: dato un insieme di veicoli e un
insieme di clienti geograficamente distribuiti, si devono progettare i percorsi
di costo minimo che servono tutti i clienti partendo da un deposito centrale e
ritornandovi. Nella variante \emph{capacitata} (CVRP), oggetto
di questo progetto, ogni veicolo ha una capacità di carico limitata e la somma
delle domande dei clienti serviti da un percorso non può superarla.

Il CVRP è \textbf{NP-hard}: contiene come caso particolare il \emph{Travelling
Salesman Problem} (basta porre la capacità pari alla domanda totale, ottenendo
un unico percorso che visita tutti i clienti). Non è quindi ragionevole
attendersi un algoritmo esatto efficiente per istanze di dimensione
interessante, e il ricorso a una metaeuristica è pienamente giustificato:
rinunciamo alla garanzia di ottimalità in cambio di soluzioni di alta qualità in
tempi accettabili.

\paragraph{Contributi del lavoro.} Il progetto sviluppa: una
\textbf{rappresentazione a giant tour} con \textbf{decodifica esatta} tramite
programmazione dinamica (Sez.~\ref{sec:split}), che garantisce l'ammissibilità per
costruzione e disaccoppia i due sotto-problemi del CVRP; un \textbf{motore
immunologico a selezione clonale} con ipermutazione inversamente proporzionale
alla \emph{fitness}, \emph{aging} elitista e selezione $(\mu+\lambda)$ con
soppressione dei duplicati (Sez.~\ref{sec:immune}); una \textbf{ricerca locale a
cinque vicinati} integrata in schema \emph{lamarckiano} con tetto di valutazioni
per invocazione (Sez.~\ref{sec:ls}); una \textbf{gestione del vincolo di flotta}
per decodifica vincolata e selezione \emph{feasibility-first}, imposta dalle
istanze a capacità satura (Sez.~\ref{sec:difficolta}); e un'\textbf{analisi critica
quantitativa} che scompone lo scarto residuo e ne individua la causa
(Sez.~\ref{sec:critica}).

\section{Il problema}
\label{sec:problema}

Sia $G=(V,E)$ un grafo completo con $V=\{v_0,v_1,\dots,v_n\}$, dove $v_0$ è il
\emph{depot} e $V\setminus\{v_0\}$ è l'insieme dei clienti. A ogni arco $(i,j)$
è associato un costo non negativo $d_{ij}$. Ogni cliente $i$ ha una domanda
$q(i)>0$; sono disponibili $m$ veicoli, ciascuno di capacità $\sigma$.

Una soluzione è un insieme di percorsi $\eta_1,\dots,\eta_m$ tale che ogni cliente
sia visitato esattamente una volta da un solo veicolo, ogni percorso inizi e
termini al depot e, per ogni veicolo $h$, valga
$\sum_{i \in V_{\eta_h}} q(i) \le \sigma$. L'obiettivo è minimizzare il costo
totale $\sum_{h=1}^{m}\sum_{(i,j)\in\eta_h} d_{ij}$.

\paragraph{Le istanze e la matrice delle distanze.}
Le dieci istanze richieste appartengono ai gruppi A, B, E e P della libreria
CVRPLIB, in formato TSPLIB con \texttt{EDGE\_WEIGHT\_TYPE = EUC\_2D}. Le
distanze sono \emph{intere}, ottenute arrotondando la distanza euclidea:
\[
  d_{ij} \;=\; \left\lfloor \sqrt{(x_i-x_j)^2 + (y_i-y_j)^2} + 0{,}5 \right\rfloor .
\]
La funzione \texttt{round} dei linguaggi più diffusi
implementa l'arrotondamento \emph{bancario} (il caso $0{,}5$ va al pari), che su
alcune coppie di punti differisce da quello TSPLIB e renderebbe i costi non
confrontabili con gli ottimi pubblicati. La matrice è stata validata
sperimentalmente: per tutte e dieci le istanze, il costo delle soluzioni ottime
ufficiali ricalcolato con la nostra matrice coincide \emph{esattamente} con quello
dichiarato da CVRPLIB --- verifica che fonda ogni confronto quantitativo del
seguito.

\paragraph{Una proprietà comune alle dieci istanze.}
Su tutte e dieci le istanze vale $\lceil \sum_i q(i) / \sigma \rceil = k$, dove
$k$ è il numero di veicoli indicato nel nome dell'istanza. Il numero di veicoli
disponibili coincide dunque con il \emph{minimo teorico}: la capacità è
\emph{satura} (il riempimento medio dei veicoli va dall'$87\%$ al $95\%$). Come
si vedrà nella Sezione~\ref{sec:difficolta}, questa caratteristica --- notata
durante l'analisi preliminare dei dati e non evidente a priori --- ha avuto
conseguenze profonde sulla progettazione dell'algoritmo.

\section{Rappresentazione e decodifica esatta (Split)}
\label{sec:split}

\subsection{La scelta della rappresentazione}

Il CVRP contiene, intrecciati, due sotto-problemi di natura diversa:
\begin{itemize}\itemsep1pt
  \item l'\textbf{assegnamento}: \emph{quali} clienti affidare a \emph{quale} veicolo (un problema di partizione con vincolo di capacità);
  \item l'\textbf{ordinamento}: in \emph{quale ordine} visitare i clienti di ciascun veicolo (un TSP su ogni percorso).
\end{itemize}

La rappresentazione più immediata --- una lista di percorsi espliciti --- obbliga
gli operatori genetici a manipolare simultaneamente entrambi gli aspetti, e
soprattutto a preoccuparsi costantemente dell'ammissibilità: qualunque mutazione
che sposti un cliente rischia di violare la capacità di un veicolo, imponendo o
meccanismi di riparazione o funzioni di penalità.

Abbiamo scelto la strada opposta. Una soluzione è codificata come un
\textbf{giant tour}: una singola permutazione $S=(s_1,\dots,s_{n})$ di tutti i
clienti, \emph{priva di separatori di percorso}. La suddivisione in percorsi non
è codificata nel genotipo: viene \emph{calcolata in modo ottimo} al momento della
valutazione, dal decodificatore \emph{Split}.

\subsection{Il decodificatore Split}

Fissato l'ordine $S$, si consideri il grafo ausiliario aciclico $H$ con nodi
$0,1,\dots,n$, in cui l'arco $(i,j)$ con $i<j$ rappresenta la decisione ``un
veicolo serve i clienti $s_{i+1},\dots,s_j$ in quest'ordine''. L'arco esiste solo
se il carico del segmento non supera la capacità, e il suo costo è
\[
  c(i,j) \;=\; d_{0,s_{i+1}} \;+\; \sum_{t=i+1}^{j-1} d_{s_t,s_{t+1}} \;+\; d_{s_j,0}.
\]
Poiché ogni cammino da $0$ a $n$ in $H$ corrisponde biunivocamente a una
suddivisione di $S$ in percorsi ammissibili, il \textbf{cammino di costo minimo}
da $0$ a $n$ fornisce la \emph{migliore} suddivisione possibile del giant tour.
$H$ è un DAG già ordinato topologicamente, quindi il cammino minimo si calcola
con la ricorrenza di Bellman (Algoritmo~\ref{alg:split}), un'applicazione diretta
della \textbf{programmazione dinamica}.

\begin{figure}[H]
\centering
\begin{tikzpicture}[scale=0.95, every node/.style={font=\small}]
  % giant tour
  \node[font=\small\bfseries] at (-1.6,1.6) {giant tour};
  \foreach \i/\x in {1/0, 2/1.1, 3/2.2, 4/3.3, 5/4.4, 6/5.5, 7/6.6} {
    \node[draw,circle,minimum size=6.5mm,fill=gray!8] (s\i) at (\x,1.6) {$s_{\i}$};
  }
  % nodi DAG
  \node[font=\small\bfseries] at (-1.6,0) {grafo $H$};
  \foreach \i/\x in {0/-0.55, 1/0.55, 2/1.65, 3/2.75, 4/3.85, 5/4.95, 6/6.05, 7/7.15} {
    \node[draw,circle,minimum size=5mm,inner sep=0pt,fill=white] (n\i) at (\x,0) {\tiny \i};
  }
  \foreach \a/\b in {0/1, 1/2, 2/3, 3/4, 4/5, 5/6, 6/7} {
    \draw[-{Latex[length=1.6mm]},gray!55] (n\a) -- (n\b);
  }
  % archi lunghi (segmenti = percorsi)
  \draw[-{Latex[length=1.8mm]},thick,blue!70!black] (n0) to[bend right=32] node[below,font=\tiny,pos=0.5] {$c(0,3)$} (n3);
  \draw[-{Latex[length=1.8mm]},thick,blue!70!black] (n3) to[bend right=32] node[below,font=\tiny,pos=0.5] {$c(3,5)$} (n5);
  \draw[-{Latex[length=1.8mm]},thick,blue!70!black] (n5) to[bend right=32] node[below,font=\tiny,pos=0.5] {$c(5,7)$} (n7);
  \draw[gray!40,dashed] (n2) to[bend left=25] (n5);
  \draw[gray!40,dashed] (n1) to[bend left=22] (n4);
  % percorsi risultanti
  \node[font=\small\bfseries] at (-1.6,-2.15) {percorsi};
  \node[draw,rounded corners,fill=blue!6,minimum height=6mm] at (0.55,-2.15) {\small $0\!\to\! s_1 s_2 s_3 \!\to\! 0$};
  \node[draw,rounded corners,fill=blue!6,minimum height=6mm] at (3.4,-2.15) {\small $0\!\to\! s_4 s_5 \!\to\! 0$};
  \node[draw,rounded corners,fill=blue!6,minimum height=6mm] at (6.2,-2.15) {\small $0\!\to\! s_6 s_7 \!\to\! 0$};
  \draw[-{Latex[length=1.6mm]},blue!70!black,thick] (2.8,-0.85) -- (2.8,-1.75);
  \node[font=\tiny,anchor=west] at (2.95,-1.3) {cammino minimo $0\!\to\! n$};
\end{tikzpicture}
\caption{Decodifica Split. Il giant tour (in alto) non contiene separatori. Nel
grafo ausiliario $H$ ogni arco rappresenta un possibile percorso (un segmento
\emph{contiguo} del tour, ammissibile rispetto alla capacità); tratteggiati alcuni
archi alternativi non scelti. Il cammino minimo da $0$ a $n$ (in blu) individua la
suddivisione ottima, da cui si leggono i percorsi.}
\label{fig:split}
\end{figure}

\begin{algorithm}[H]
\SetAlgoLined
\DontPrintSemicolon
\KwIn{giant tour $S$, domande $q$, distanze $d$, capacità $\sigma$}
\KwOut{costo ottimo $V[n]$ e suddivisione in percorsi}
$V[0] \leftarrow 0$; \quad $V[j] \leftarrow +\infty$, \; $\mathit{pred}[j] \leftarrow -1$ \quad per $j=1,\dots,n$\;
\For{$i \leftarrow 0$ \KwTo $n-1$}{
  \lIf{$V[i] = +\infty$}{\textbf{continue}}
  $\mathit{load} \leftarrow 0$; \; $\mathit{cost} \leftarrow 0$\;
  \For{$j \leftarrow i+1$ \KwTo $n$}{
    $\mathit{load} \leftarrow \mathit{load} + q(s_j)$\;
    \lIf{$\mathit{load} > \sigma$}{\textbf{break} \quad \tcp*[h]{potatura: nessuna estensione è ammissibile}}
    \eIf{$j = i+1$}{
      $\mathit{cost} \leftarrow d_{0,s_j} + d_{s_j,0}$ \tcp*{percorso con un solo cliente}
    }{
      $\mathit{cost} \leftarrow \mathit{cost} - d_{s_{j-1},0} + d_{s_{j-1},s_j} + d_{s_j,0}$ \tcp*{estensione in $O(1)$}
    }
    \If{$V[i] + \mathit{cost} < V[j]$}{
      $V[j] \leftarrow V[i] + \mathit{cost}$; \quad $\mathit{pred}[j] \leftarrow i$\;
    }
  }
}
\Return $V[n]$ e i percorsi ottenuti risalendo $\mathit{pred}$ da $n$ a $0$\;
\caption{\textsc{Split} --- suddivisione ottima di un giant tour (Bellman su DAG)}
\label{alg:split}
\end{algorithm}

\paragraph{Complessità.} Il ciclo interno si interrompe non appena il carico
supera $\sigma$: detto $s^{*}$ il numero massimo di clienti che stanno in un
veicolo, il costo è $O(n \cdot s^{*})$. L'aggiornamento incrementale di
$\mathit{cost}$ in $O(1)$ (riga~9) evita di ricalcolare da capo il costo di ogni
segmento. In pratica, su $n=100$ una decodifica richiede meno di un millisecondo.

\paragraph{Perché questa scelta paga.} Quattro vantaggi concreti.
\emph{Ammissibilità per costruzione}: ogni soluzione decodificata rispetta la
capacità, senza riparazioni né penalità. \emph{Operatori semplici}: le mutazioni
agiscono su una permutazione e non possono produrre soluzioni inammissibili.
\emph{Ottimalità della suddivisione}: a parità di ordine nessun algoritmo può fare
meglio della programmazione dinamica, e la metaeuristica può concentrarsi
\emph{solo} sull'ordinamento. \emph{Paesaggio più regolare}: una singola mutazione
può riorganizzare radicalmente i percorsi, perché la decodifica si riadatta.

\paragraph{Verifica di correttezza.} Oltre al confronto con un'enumerazione
esaustiva su istanze giocattolo, il decodificatore soddisfa una proprietà
verificabile in modo esatto: concatenando i percorsi della soluzione
\emph{ottima} di CVRPLIB in un giant tour e applicandovi lo Split (vincolato a
$k$ percorsi) si riottiene \emph{esattamente} il costo ottimo. La dimostrazione è
immediata: la suddivisione ottima è una delle suddivisioni contigue con al più
$k$ percorsi, dunque il costo calcolato non può essere superiore all'ottimo; e
ogni suddivisione con al più $k$ percorsi è una soluzione ammissibile del CVRP,
dunque non può essere inferiore. Questo controllo è superato su tutte e dieci le
istanze.

\section{L'algoritmo immunologico a selezione clonale}
\label{sec:immune}

\subsection{Motivazione della scelta}

Fra le cinque opzioni proposte, la scelta è caduta sull'\textbf{algoritmo
immunologico a selezione clonale} (opzione~2), ibridato con ricerca locale. Le
ragioni sono tre.

In primo luogo, il paradigma immunologico offre un \textbf{meccanismo di
bilanciamento fra esplorazione e sfruttamento che è intrinseco e adattivo}, e non
un parametro fisso: l'ipermutazione modula l'intensità della perturbazione in
funzione della qualità dell'individuo, cosicché le soluzioni migliori vengono
raffinate con piccoli aggiustamenti mentre quelle peggiori esplorano
aggressivamente lo spazio. In un algoritmo genetico classico il tasso di
mutazione è invece uniforme sulla popolazione. In secondo luogo, l'\emph{aging}
fornisce un meccanismo esplicito e diretto contro la convergenza prematura,
problema notoriamente critico nel CVRP, dove il paesaggio di ricerca è ricco di
ottimi locali profondi. In terzo luogo, la selezione clonale non richiede un
operatore di ricombinazione: il crossover su permutazioni (OX, PMX, ecc.) è
sempre un compromesso delicato, e poterne fare a meno riduce le scelte arbitrarie
e le fonti di errore.



\subsection{Struttura di un anticorpo e ciclo generazionale}

Un \textbf{anticorpo} è la tripla $(\text{giant tour},\ \text{costo},\ \text{età})$.
Il costo (la \emph{fitness}, da minimizzare) è ottenuto dalla decodifica della
Sezione~\ref{sec:split}; l'età governa l'invecchiamento. La popolazione ha
dimensione $B$ e il ciclo generazionale è riassunto nella
Figura~\ref{fig:pipeline} e nell'Algoritmo~\ref{alg:immune}.

\begin{figure}[H]
\centering
\begin{tikzpicture}[
  box/.style={draw,rounded corners,minimum height=7.5mm,minimum width=21mm,
              align=center,font=\small,fill=gray!10},
  op/.style={draw,rounded corners,minimum height=7.5mm,minimum width=21mm,
             align=center,font=\small,fill=orange!10},
  dec/.style={draw,rounded corners,minimum height=7.5mm,minimum width=19mm,
              align=center,font=\small,fill=blue!8},
  arr/.style={-{Latex[length=2mm]},thick}
]
  % riga superiore
  \node[box]                          (init)  at (0,0)    {inizializzazione\\{\scriptsize greedy + casuale}};
  \node[op, right=6mm of init]        (clone) {\textbf{cloning}\\{\scriptsize $\times\,\mathit{dup}$}};
  \node[op, right=6mm of clone]       (hyper) {\textbf{ipermutazione}\\{\scriptsize $\propto$ inversa fitness}};
  \node[op, right=6mm of hyper]       (eval)  {\textbf{valutazione}\\{\scriptsize Split (DP)}};
  \node[op, right=6mm of eval]        (ls)    {\textbf{ricerca locale}\\{\scriptsize sui $K$ migliori}};
  % riga inferiore
  \node[op, below=9mm of ls]          (aging) {\textbf{aging}\\{\scriptsize elitista}};
  \node[op, left=6mm of aging]        (sel)   {\textbf{selezione} $(\mu\!+\!\lambda)$\\{\scriptsize senza duplicati}};
  \node[dec, left=6mm of sel]         (stop)  {budget FE\\esaurito?};
  \node[box, left=6mm of stop]        (out)   {migliore\\soluzione};

  \draw[arr] (init)  -- (clone);
  \draw[arr] (clone) -- (hyper);
  \draw[arr] (hyper) -- (eval);
  \draw[arr] (eval)  -- (ls);
  \draw[arr] (ls)    -- (aging);
  \draw[arr] (aging) -- (sel);
  \draw[arr] (sel)   -- (stop);
  \draw[arr] (stop)  -- node[above,font=\scriptsize] {sì} (out);
  % retroazione: nuova generazione -> riparte dal cloning (NON dall'inizializzazione)
  \draw[arr] (stop.south) -- ++(0,-5mm) -| node[pos=0.25,below,font=\scriptsize]
        {no: nuova generazione} (clone.south);
\end{tikzpicture}
\caption{Ciclo generazionale. L'inizializzazione avviene una sola volta; le fasi
in arancione sono gli operatori immunologici e l'ibridazione con la ricerca
locale, ripetuti a ogni generazione finché il budget di valutazioni non è
esaurito.}
\label{fig:pipeline}
\end{figure}

\paragraph{Inizializzazione: greedy e casuale.}
Il $25\%$ della popolazione è seminato a partire dall'euristica costruttiva di
\textbf{Clarke--Wright} (algoritmo dei \emph{savings}, un tipico algoritmo
\emph{greedy}): si parte dalla soluzione banale con un veicolo per cliente e si
fondono iterativamente i percorsi in ordine decrescente di risparmio
$s(i,j)=d_{0i}+d_{0j}-d_{ij}$, quando ammissibile. Il primo individuo è la
soluzione di Clarke--Wright pura; gli altri ne sono perturbazioni casuali. Il
restante $75\%$ è costituito da permutazioni casuali.

Il bilanciamento è un compromesso: una popolazione interamente greedy sarebbe fatta
di individui quasi identici (diversità nulla), una interamente casuale sprecherebbe
gran parte del budget per raggiungere una qualità che l'euristica costruttiva
fornisce gratuitamente.

\subsection{Ipermutazione inversamente proporzionale alla fitness}

È l'operatore caratterizzante del paradigma immunologico. Detti $c_i$ il costo
dell'$i$-esimo anticorpo, $c_{\min}$ e $c_{\max}$ il costo minimo e massimo nella
popolazione corrente, si definisce la \textbf{fitness normalizzata}
\[
  \hat{f}_i \;=\; \frac{c_{\max}-c_i}{c_{\max}-c_{\min}} \;\in\; [0,1],
  \qquad (\hat{f}=1 \text{ per il migliore},\; \hat{f}=0 \text{ per il peggiore})
\]
e il \textbf{numero di mutazioni elementari} da applicare al clone:
\[
  \boxed{\;M_i \;=\; \max\!\left(1,\; \left\lceil\, n \cdot e^{-\rho\,\hat{f}_i} \,\right\rceil\right)\;}
  \qquad\text{con } \rho = \ln n .
\]
La scelta $\rho=\ln n$ non è arbitraria: con essa il migliore riceve
\emph{esattamente una} mutazione ($n \cdot e^{-\ln n} = 1$) e il peggiore ne
riceve $n$, cioè una randomizzazione pressoché completa. La legge esponenziale
realizza dunque, con un solo parametro, l'intero spettro fra sfruttamento puro e
esplorazione pura, assegnando a ciascun individuo l'intensità di perturbazione
adatta alla sua qualità. Nel caso degenere $c_{\max}=c_{\min}$ si pone
$\hat{f}=1$ per tutti; la diversità viene allora reintrodotta dall'aging e dalla
soppressione dei duplicati, non da mutazioni più violente.

Ogni mutazione elementare è estratta fra tre operatori: \textbf{inversione} di un
segmento (peso $0{,}50$), \textbf{scambio} di due posizioni ($0{,}25$) e
\textbf{inserimento} di un cliente in altra posizione ($0{,}25$). All'inversione è
assegnato il peso maggiore perché è la mossa più efficace sulle permutazioni di
percorso: preserva tutte le adiacenze interne al segmento, distruggendo solo i due
archi agli estremi. Tutte e tre preservano l'invariante di permutazione;
l'ammissibilità della soluzione decodificata è poi garantita dallo Split.

\subsection{Aging elitista e selezione}

\paragraph{Aging.} Ogni anticorpo ha un'età, incrementata a ogni generazione. Un
clone \emph{eredita} l'età del genitore, ma essa viene \textbf{azzerata se il
clone migliora il genitore}: le linee evolutive che stanno progredendo vengono
così protette, mentre quelle stagnanti invecchiano e muoiono. Superata l'età
massima $\tau_B$, l'anticorpo viene rimosso. L'aging è \textbf{elitista}: il
miglior anticorpo corrente non muore mai, per non perdere l'informazione
migliore disponibile.

\paragraph{Selezione $(\mu+\lambda)$ con soppressione dei duplicati.} Dall'unione
di genitori e cloni sopravvissuti all'aging si selezionano i $B$ migliori. La
selezione, però, \textbf{scarta i duplicati}: a parità di costo passa un solo
anticorpo, e i duplicati rientrano solo se non ci sono abbastanza soluzioni
distinte. Questa regola, apparentemente minore, si è rivelata essenziale
(Sezione~\ref{sec:difficolta}). Se dopo la selezione la popolazione è sotto
organico, si generano \textbf{nuovi anticorpi casuali} (nascite), ulteriore
iniezione di diversità.

\begin{algorithm}[H]
\SetAlgoLined
\DontPrintSemicolon
\KwIn{istanza, budget $\mathrm{FE}_{\max}=3.5\times10^5$, parametri $B$, $\mathit{dup}$, $\rho$, $\tau_B$, $K$, $\mathit{cap}$}
\KwOut{migliore soluzione ammissibile trovata}
$P \leftarrow$ inizializza $B$ anticorpi ($25\%$ da Clarke--Wright, resto casuali); valuta\;
\While{budget non esaurito}{
  $\hat{f} \leftarrow$ fitness normalizzata di $P$\;
  $C \leftarrow \emptyset$\;
  \ForEach{$a \in P$}{
    $M \leftarrow \max(1, \lceil n\, e^{-\rho \hat{f}_a}\rceil)$ \tcp*{ipermutazione}
    \For{$1$ \KwTo $\mathit{dup}$}{
      $c \leftarrow$ \textsc{Muta}($a$, $M$); \; valuta $c$ con \textsc{Split} \tcp*{1 FE}
      $c.\mathit{et\grave{a}} \leftarrow a.\mathit{et\grave{a}}$; \quad $C \leftarrow C \cup \{c\}$\;
    }
  }
  \ForEach{$c \in$ i $K$ migliori di $C$}{
    $c \leftarrow$ \textsc{RicercaLocale}($c$, tetto $=\mathit{cap}$ FE) \tcp*{ibridazione lamarckiana}
  }
  \lForEach{$c \in C$ che migliora il proprio genitore}{$c.\mathit{et\grave{a}} \leftarrow 0$}
  $\text{incrementa l'età di } P \cup C$; \; rimuovi chi ha età $> \tau_B$ (mai il migliore) \tcp*{aging}
  $P \leftarrow$ i $B$ migliori dei sopravvissuti, \textbf{senza duplicati} \tcp*{selezione $(\mu+\lambda)$}
  \lIf{$|P| < B$}{completa con anticorpi casuali (nascite)}
}
\caption{Algoritmo immunologico a selezione clonale, ibridato con ricerca locale}
\label{alg:immune}
\end{algorithm}

\section{Ricerca locale e ibridazione}
\label{sec:ls}

L'ibridazione con la ricerca locale è ciò che porta la qualità delle soluzioni
vicino agli ottimi noti: è l'aspetto \emph{hybrid metaheuristics} del progetto.
La ricerca locale opera sui \emph{percorsi decodificati} e usa cinque vicinati,
esplorati in sequenza con strategia \emph{first-improvement} finché un giro
completo non produce miglioramenti (Figura~\ref{fig:mosse}): \textbf{2-opt}
(intra-percorso: inverte un segmento, eliminando gli incroci); \textbf{Or-opt}
(sposta una catena di $1$--$3$ clienti consecutivi, eventualmente invertendola,
nello stesso percorso o in un altro); \textbf{relocate} (sposta un cliente in un
altro percorso: è il caso di Or-opt con catena unitaria); \textbf{swap} (scambia
due clienti di percorsi diversi); \textbf{2-opt*} (scambia le \emph{code} di due
percorsi, e come caso limite ne realizza la \emph{fusione}). I primi due agiscono
sull'\emph{ordinamento}, gli ultimi tre sull'\emph{assegnamento}: insieme coprono
entrambe le componenti del problema.

\begin{figure}[H]
\centering
\begin{tikzpicture}[scale=0.78, every node/.style={font=\scriptsize},
    dep/.style={draw,fill=black,rectangle,minimum size=2.2mm,inner sep=0pt},
    cli/.style={draw,fill=gray!35,circle,minimum size=2.2mm,inner sep=0pt},
    old/.style={gray!50,dashed,thick},
    new/.style={blue!70!black,thick}]
  % 2-opt
  \begin{scope}[shift={(0,0)}]
    \node[font=\small\bfseries] at (1.6,2.1) {2-opt};
    \node[dep] (d) at (0,0) {};
    \node[cli] (a) at (1,1.2) {}; \node[cli] (b) at (2.2,1.3) {};
    \node[cli] (c) at (2.6,0.2) {}; \node[cli] (e) at (1.3,0.1) {};
    \draw[old] (d) -- (a) -- (c) -- (b) -- (e) -- (d);
    \draw[new] (d) -- (a) -- (b) -- (c) -- (e) -- (d);
    \node at (1.6,-0.75) {inverte un segmento};
  \end{scope}
  % relocate / or-opt
  \begin{scope}[shift={(4.4,0)}]
    \node[font=\small\bfseries] at (1.6,2.1) {Or-opt / relocate};
    \node[dep] (d2) at (0,0) {};
    \node[cli] (a2) at (0.9,1.3) {}; \node[cli] (b2) at (2.1,1.4) {};
    \node[cli] (c2) at (2.7,0.3) {}; \node[cli,fill=orange!60] (x2) at (1.5,0.75) {};
    \draw[old] (a2) -- (x2) -- (b2);
    \draw[new] (d2) -- (a2) -- (b2) -- (c2) -- (d2);
    \draw[new] (c2) to[bend right=20] (x2); \draw[new] (x2) to[bend right=25] (d2);
    \node at (1.6,-0.75) {sposta una catena};
  \end{scope}
  % swap
  \begin{scope}[shift={(8.8,0)}]
    \node[font=\small\bfseries] at (1.6,2.1) {swap};
    \node[dep] (d3) at (0,0) {};
    \node[cli] (a3) at (0.8,1.35) {}; \node[cli,fill=orange!60] (u3) at (1.9,1.5) {};
    \node[cli] (b3) at (2.9,0.9) {}; \node[cli,fill=green!45] (v3) at (2.0,0.35) {};
    \node[cli] (c3) at (0.9,0.4) {};
    \draw[new] (d3) -- (a3) -- (u3) -- (b3);
    \draw[new] (d3) -- (c3) -- (v3);
    \draw[<->,orange!80!black,thick,dashed] (u3) -- (v3);
    \node at (1.6,-0.75) {scambia due clienti};
  \end{scope}
  % 2-opt*
  \begin{scope}[shift={(13.2,0)}]
    \node[font=\small\bfseries] at (1.4,2.1) {2-opt*};
    \node[dep] (d4) at (0,0) {};
    \node[cli] (p1) at (0.9,1.4) {}; \node[cli] (p2) at (2.2,1.5) {};
    \node[cli] (q1) at (1.0,0.5) {}; \node[cli] (q2) at (2.3,0.5) {};
    \draw[old] (p1) -- (p2); \draw[old] (q1) -- (q2);
    \draw[new] (d4) -- (p1) -- (q2); \draw[new] (d4) -- (q1) -- (p2);
    \node at (1.4,-0.75) {scambia le code};
  \end{scope}
\end{tikzpicture}
\caption{I vicinati della ricerca locale (tratteggiato grigio: archi rimossi; blu:
archi creati). I primi due agiscono sull'ordinamento interno ai percorsi, gli
ultimi due sull'assegnamento dei clienti ai veicoli.}
\label{fig:mosse}
\end{figure}

\paragraph{Valutazione incrementale e liste di candidati.} Ogni mossa è valutata
in $O(1)$ tramite la variazione di costo $\Delta$ (differenza fra archi creati e
archi rimossi); i controlli di capacità sono anch'essi in $O(1)$ grazie alle
somme prefisse dei carichi. L'esplorazione completa dei vicinati
inter-percorso costerebbe $O(n^2)$ valutazioni per passata, insostenibile con un
budget di $3{,}5\times10^5$ valutazioni: si adotta quindi la tecnica standard
delle \textbf{liste di candidati}, considerando per ogni cliente solo mosse che
creano archi verso uno dei suoi $h=20$ vicini più prossimi.

\paragraph{Ibridazione lamarckiana.} La soluzione migliorata dalla ricerca locale
viene \textbf{riscritta nel genotipo}: i percorsi ottimizzati sono riconcatenati
in un nuovo giant tour, che sostituisce quello originale dell'anticorpo. Si tratta
di apprendimento \emph{lamarckiano} (il carattere acquisito viene ereditato), in
contrapposizione allo schema \emph{baldwiniano}, in cui il miglioramento
influenzerebbe solo la fitness lasciando intatto il genotipo. La scelta
lamarckiana è motivata dalla struttura della rappresentazione: poiché lo Split
estrae segmenti contigui, la riconcatenazione dei percorsi ottimizzati è una
operazione esatta e senza perdita di informazione, e trasmette ai discendenti il
lavoro già svolto dalla ricerca locale.

\paragraph{Il tetto di valutazioni per invocazione.} Una discesa completa fino
all'ottimo locale costa in media $\sim\!12\,000$ valutazioni: con l'intero budget se
ne potrebbero fare solo $29$. Ma il costo non è distribuito uniformemente: gran
parte delle valutazioni è consumata dall'ultima passata, quella che scandaglia tutti
i vicinati \emph{senza trovare nulla}, per certificare l'ottimalità locale. Abbiamo
quindi introdotto un \textbf{tetto di FE per invocazione} ($\mathit{cap}$): la
discesa si interrompe al tetto, rinunciando alla certificazione ma conservando i
miglioramenti ad alto rendimento. Con $K$, è il parametro più influente
dell'algoritmo (Sezione~\ref{sec:tuning}).

\subsection{Complessità computazionale}
\label{sec:complessita}

Detti $n$ i clienti, $B$ la popolazione, $\mathit{dup}$ i cloni per anticorpo,
$s^{*}$ il massimo numero di clienti in un veicolo, $h$ l'ampiezza delle liste di
candidati, una generazione costa: $O(B\,\mathit{dup}\,n)$ per l'ipermutazione nel
caso peggiore (il peggior anticorpo riceve $n$ mutazioni, ma in media $M_i \ll n$);
$O(B\,\mathit{dup}\cdot k\,n\,s^{*})$ per la valutazione dei cloni (una decodifica
ciascuno, con la DP vincolata a $k$ percorsi); $O(K\cdot\mathit{cap})$ per la
ricerca locale, limitata per costruzione dal tetto, con una passata sui vicinati
inter-percorso che costa $O(n\,h)$ anziché $O(n^2)$ grazie alle liste di candidati;
e $O(B\,\mathit{dup}\log(B\,\mathit{dup}))$ per aging e selezione.

Con i parametri congelati il termine dominante è di gran lunga la ricerca locale
($K\cdot\mathit{cap}=15\,000$ valutazioni contro le $B\cdot\mathit{dup}=45$
decodifiche): è questa asimmetria a spiegare perché il budget di $3{,}5\times10^5$
valutazioni si esaurisca in circa $24$ generazioni, e perché il tetto
$\mathit{cap}$ sia il parametro più influente dell'algoritmo. In termini di tempo,
una \emph{run} completa richiede fra $1{,}5$ e $6$ secondi.

\paragraph{Convenzione di conteggio delle valutazioni.} Il criterio di arresto
imposto è il numero massimo di valutazioni della funzione di fitness. Adottiamo
la convenzione più \textbf{conservativa}: \emph{una valutazione $=$ il calcolo del
costo di una soluzione candidata}, sia essa una decodifica Split completa
\emph{oppure} la valutazione incrementale $\Delta$ di un vicino esaminato dalla
ricerca locale. Il fatto che il $\Delta$ si calcoli in $O(1)$ è un'ottimizzazione
di \emph{velocità}, non una ragione per non contare la valutazione. Una
convenzione più permissiva (contare solo le decodifiche complete) consentirebbe
una ricerca locale molto più estesa a parità di budget dichiarato; la nostra
scelta va nella direzione opposta e non può in alcun caso sovrastimare le
prestazioni. L'arresto è esatto: tutte le $50$ esecuzioni terminano a
$350\,000$ valutazioni, né una in più né una in meno.

\section{Difficoltà affrontate e scelte progettuali}
\label{sec:difficolta}

Questa sezione documenta i tre problemi che hanno richiesto le decisioni
progettuali più delicate. Sono presentati nella forma in cui si sono
effettivamente manifestati (sintomo, diagnosi, rimedio).

\subsection{Convergenza prematura: la soppressione dei duplicati}

\textbf{Sintomo.} Nelle prime esecuzioni complete, l'ultimo miglioramento della
soluzione migliore avveniva sistematicamente intorno alla generazione~14, dopo la
quale l'algoritmo consumava il restante $75\%$ del budget senza progressi.

\textbf{Diagnosi.} La selezione $(\mu+\lambda)$, prendendo semplicemente i $B$
migliori, riempiva la popolazione di \emph{copie} dello stesso ottimo locale.
Ma se tutti gli anticorpi hanno lo stesso costo, allora $c_{\max}=c_{\min}$, la
fitness normalizzata degenera ($\hat{f}=1$ per tutti) e l'ipermutazione si
riduce al minimo (una sola mutazione) \emph{per tutti}: l'algoritmo perde
esattamente il meccanismo che dovrebbe salvarlo dalla stagnazione. La legge di
ipermutazione, da sola, non basta: presuppone diversità nella popolazione.

\textbf{Rimedio.} La \emph{soppressione dei duplicati} in selezione: a parità di
costo passa un solo anticorpo. L'effetto è stato immediato e misurabile: su
A-n45-k7 lo scarto è sceso dal $2{,}97\%$ all'$1{,}66\%$, e soprattutto
\emph{l'ultimo miglioramento si è spostato dalla generazione~14 alle
generazioni~42--54}, cioè l'algoritmo ha ripreso a progredire fino alla fine del
budget.

\subsection{Il vincolo di flotta e le istanze a capacità satura}

\textbf{Sintomo.} Sull'istanza P-n50-k10 --- riempimento medio dei veicoli al
$95\%$ --- l'algoritmo non produceva \emph{alcuna} soluzione con al più $k=10$
percorsi.

\textbf{Diagnosi.} Lo Split minimizza il \emph{costo}, non il \emph{numero di
percorsi}. Su istanze a capacità satura, la suddivisione a costo minimo di un
giant tour usa quasi sempre più di $k$ percorsi: spezzare un percorso in due
allunga il costo di poco (due brevi tratte in più verso il depot) ma rilassa
enormemente il vincolo di capacità. La misura empirica è impressionante: la
percentuale di cloni la cui decodifica \emph{non ammette} una suddivisione in
$\le k$ percorsi arriva al $93\%$ su P-n50-k10 e all'$84\%$ su B-n66-k9
(Tabella~\ref{tab:operatori}). Il vincolo sul numero di veicoli, che sulle
istanze ``facili'' è automaticamente soddisfatto, qui è il vincolo dominante.

\textbf{Rimedio, in due passi.} Il primo tentativo --- decodificare normalmente e
poi, solo alla fine, forzare una suddivisione vincolata sulla migliore soluzione
--- ha prodotto soluzioni conformi ma pessime ($16{,}5\%$ di scarto): la ricerca
non stava affatto ottimizzando ciò che poi veniva riportato. Il secondo tentativo
--- usare direttamente la \textbf{decodifica vincolata} a $k$ percorsi come
funzione di fitness (una programmazione dinamica bidimensionale
$V[r][j]$, con $r$ = numero di percorsi usati) --- ha peggiorato ulteriormente
($34\%$): i costi della decodifica \emph{rilassata} (senza vincolo sul numero di
percorsi), essendo per costruzione più bassi, dominavano in selezione, e gli
anticorpi conformi venivano sistematicamente eliminati.

La soluzione corretta richiede \emph{entrambi} gli ingredienti: decodifica
vincolata come fitness \textbf{e} selezione \textbf{lessicografica
feasibility-first}, in cui un anticorpo conforme (al più $k$ percorsi) precede
\emph{sempre} un anticorpo non conforme, e solo a parità di conformità decide il
costo. Con questa regola di dominanza lo scarto su P-n50-k10 è crollato dal
$34\%$ allo $0{,}57\%$. Si noti che non si tratta di un approccio a
\emph{penalty function} (opzione~5 della consegna): non c'è alcun coefficiente di
penalità da calibrare, né una esplorazione controllata della regione non
ammissibile; è una semplice regola di dominanza in selezione, che mantiene le
soluzioni non conformi nella popolazione come materiale genetico ma impedisce
loro di prevalere.

\subsection{L'ipermutazione troppo aggressiva}

Un tentativo di aumentare l'esplorazione dimezzando $\rho$ (a $\ln(n)/2$, così
che anche i migliori ricevano più mutazioni) ha prodotto un risultato
paradossale: la soluzione migliore riportata restava bloccata alla \emph{prima}
valutazione. La spiegazione è ora evidente alla luce del punto precedente:
mutazioni più violente producono tour che si decodificano in più di $k$ percorsi,
i quali non possono diventare la soluzione riportata. La lezione è che i
parametri di un algoritmo non sono indipendenti dai vincoli del problema: $\rho =
\ln n$ non è solo elegante, è anche l'unico valore che mantiene i cloni
all'interno della regione ammissibile.

\section{Taratura dei parametri}
\label{sec:tuning}

La taratura è stata condotta con una strategia di \emph{coordinate descent} in
più round, su un sottoinsieme di quattro istanze eterogenee (una per gruppo:
A-n60-k9, B-n66-k9, E-n101-k14, P-n101-k4) e con \textbf{tre semi casuali
dedicati (100--102), disgiunti dai semi $0$--$4$ usati per gli esperimenti
finali}. Questa separazione è una precauzione metodologica essenziale: tarare i
parametri sugli stessi campioni casuali con cui poi si misurano le prestazioni
significherebbe misurare sui dati di addestramento, sovrastimando la qualità.

\begin{table}[H]
\centering
\small
\caption{Taratura dei parametri: scarto percentuale medio dall'ottimo sulle
$4$ istanze di taratura $\times$ $3$ semi. \emph{Round~1}: forma del budget di
ricerca locale ($K$ cloni, tetto \emph{cap} di FE per invocazione).
\emph{Round~2}: struttura della popolazione e aging. \emph{Round~3}:
ri-validazione delle configurazioni di frontiera sul motore definitivo.}
\label{tab:tuning}
\begin{tabular}{llr@{\hskip 14pt}llr@{\hskip 14pt}llr}
\toprule
\multicolumn{3}{c}{\textbf{Round 1} --- budget LS} &
\multicolumn{3}{c}{\textbf{Round 2} --- popolazione, aging} &
\multicolumn{3}{c}{\textbf{Round 3} --- ri-validazione} \\
\cmidrule(r){1-3}\cmidrule(r){4-6}\cmidrule(r){7-9}
config. & & scarto & config. & & scarto & config. & & scarto \\
\midrule
$K$=3, cap=5000 & & \textbf{3,11\%} & \emph{dup}=3      & & \textbf{3,00\%} & $K$=3, cap=5000, \emph{dup}=3, $\tau_B$=20 & & \textbf{2,90\%} \\
$K$=2, cap=5000 & & 3,26\%          & $B$=10            & & 3,10\%          & $K$=2, cap=5000, \emph{dup}=3, $\tau_B$=20 & & 3,00\% \\
$K$=3, cap=3000 & & 4,01\%          & $B$=20            & & 3,19\%          & $K$=3, cap=5000, \emph{dup}=3, $\tau_B$=10 & & 3,27\% \\
$K$=3, cap=1500 & & 5,81\%          & $\tau_B$=10       & & 3,24\%          & $K$=2, cap=5000, \emph{dup}=3, $\tau_B$=10 & & 3,40\% \\
$K$=3, cap=2000 & & 7,07\%          & $\tau_B$=30       & & 3,26\%          & & & \\
$K$=4, cap=1500 & & 7,81\%          & base              & & 3,26\%          & & & \\
\bottomrule
\end{tabular}
\end{table}

I risultati (Tabella~\ref{tab:tuning}) indicano con chiarezza che \textbf{poche
discese profonde battono molte discese superficiali}: a parità di budget di
ricerca locale per generazione, le configurazioni con tetto elevato ($5000$~FE)
dominano nettamente quelle con tetto basso. L'interpretazione è che una discesa
troncata troppo presto non riesce a uscire dai vicinati più economici (2-opt e
Or-opt) e non arriva mai a esplorare le mosse inter-percorso, che sono quelle
che riorganizzano davvero la soluzione. Il Round~2 mostra poi che aumentare il
numero di cloni per genitore (\emph{dup}~$=3$) intensifica l'esplorazione attorno
alle soluzioni élite, con il beneficio maggiore proprio sull'istanza più ostica
(P-n101-k4).

\paragraph{Una nota metodologica onesta.} I Round~1 e~2 hanno preceduto le modifiche
della Sezione~\ref{sec:difficolta} (decodifica vincolata e selezione
feasibility-first). Poiché quelle modifiche cambiano l'algoritmo in modo
sostanziale, la taratura precedente non era più necessariamente valida: si è quindi
eseguito un \textbf{round di ri-validazione} sul motore definitivo, sulle quattro
configurazioni di frontiera. Ne è emerso un fatto istruttivo: $\tau_B=10$, che sul
motore intermedio era vantaggioso, sul definitivo peggiora sistematicamente ---
con la selezione feasibility-first le linee conformi impiegano più generazioni a
maturare, e un aging aggressivo le elimina prematuramente.

\paragraph{Configurazione finale (congelata).}
$B=15$ (popolazione), $\mathit{dup}=3$ (cloni per anticorpo), $\rho=\ln n$
(ipermutazione), $\tau_B=20$ (età massima), $K=3$ (cloni sottoposti a ricerca
locale), $\mathit{cap}=5000$ (tetto di FE per discesa), $h=20$ (ampiezza liste di
candidati), $25\%$ di semina greedy.

\section{Originalità introdotte}
\label{sec:originalita}

Riassumiamo qui gli elementi che riteniamo costituiscano il contributo originale
del progetto, al di là dell'applicazione dello schema standard di selezione
clonale.

\begin{enumerate}\itemsep1pt
  \item \textbf{Decodifica esatta dentro una metaeuristica.} L'uso di un algoritmo
  \emph{esatto} (programmazione dinamica) come funzione di valutazione di un
  algoritmo \emph{approssimato}: la parte del problema che sappiamo risolvere
  esattamente (la suddivisione, dato l'ordine) viene risolta esattamente; alla
  metaeuristica resta solo la parte che richiede ricerca.

  \item \textbf{Gestione del vincolo di flotta per dominanza.} Decodifica vincolata
  e selezione lessicografica \emph{feasibility-first} sono, per quanto ne sappiamo,
  una soluzione non standard al problema del numero di veicoli nella
  rappresentazione a giant tour, ed evitano del tutto la calibrazione --- notoriamente
  fragile --- di coefficienti di penalità.

  \item \textbf{Aging con azzeramento dell'età sui miglioramenti.} Protegge le linee
  evolutive produttive invece di invecchiarle meccanicamente: l'aging diventa un
  ricambio dei \emph{fallimenti} più che un ricambio anagrafico.

  \item \textbf{Soppressione dei duplicati come pre-requisito dell'ipermutazione.}
  Il legame fra collasso della popolazione e degenerazione della legge di
  ipermutazione (Sez.~\ref{sec:difficolta}) non è un accorgimento implementativo:
  è la constatazione che, in un algoritmo immunologico, mantenere la \emph{varianza}
  della fitness è condizione necessaria perché l'operatore caratterizzante del
  paradigma funzioni.

  \item \textbf{Ricerca locale a budget limitato per invocazione.} Il tetto di FE
  per discesa --- motivato dall'osservazione che la certificazione di ottimalità
  locale consuma gran parte del costo senza produrre miglioramenti --- si è rivelato
  il parametro più influente dell'intero algoritmo.
\end{enumerate}

\section{Risultati sperimentali}
\label{sec:risultati}

\paragraph{Protocollo.} Dieci istanze CVRPLIB, $5$ esecuzioni indipendenti per
istanza (semi $0$--$4$, fissati e dichiarati), criterio di arresto
$\mathrm{FE}=3{,}5\times10^5$, parametri congelati dalla taratura. L'intera
campagna ($50$ esecuzioni) richiede circa $150$ secondi su un normale portatile.
Ogni esecuzione è \textbf{riproducibile bit a bit}: a parità di seme il risultato
è identico.

\begin{table}[H]
\centering
\small
\caption{Risultati su $5$ esecuzioni indipendenti per istanza
($\mathrm{FE}=3{,}5\times10^5$). \emph{best}: miglior costo fra le $5$ esecuzioni;
\emph{mean} e \emph{dev.\ std}: media e deviazione standard campionaria dei
migliori costi di ciascuna esecuzione; \emph{iter.\ medie}: numero medio di
generazioni necessarie a raggiungere il migliore risultato di ciascuna esecuzione;
\emph{scarto}: distanza percentuale del \emph{best} dall'ottimo noto (BKS).
Tutte le città sono servite in tutte le $50$ esecuzioni (\emph{satisfability}
$=100\%$) e il numero di percorsi è sempre pari a $k$.}
\label{tab:risultati}
\begin{tabular}{lrrrrrrr}
\toprule
Istanza & BKS & best & scarto & mean & dev.\ std & iter.\ medie & percorsi/$k$ \\
\midrule
A-n45-k7   & 1146 & \textbf{1148} & 0,17\% & 1158,2 & 9,2 & 17,4 & 7/7 \\
A-n60-k9   & 1354 & \textbf{1358} & 0,30\% & 1363,8 & 4,8 & 15,4 & 9/9 \\
A-n80-k10  & 1763 & 1811 & 2,72\% & 1818,4 & 8,8 & 14,8 & 10/10 \\
B-n56-k7   & 707  & \textbf{710}  & 0,42\% & 712,6  & 1,9 & 15,6 & 7/7 \\
B-n66-k9   & 1316 & 1332 & 1,22\% & 1335,6 & 3,0 & 17,0 & 9/9 \\
B-n78-k10  & 1221 & 1239 & 1,47\% & 1245,8 & 5,1 & 10,8 & 10/10 \\
E-n76-k8   & 735  & 743  & 1,09\% & 757,6  & 8,8 & 15,2 & 8/8 \\
E-n101-k14 & 1067 & 1109 & 3,94\% & 1114,2 & 4,9 & 15,0 & 14/14 \\
P-n50-k10  & 696  & 705  & 1,29\% & 708,6  & 3,3 & 19,2 & 10/10 \\
P-n101-k4  & 681  & 719  & 5,58\% & 724,4  & 4,4 & 17,8 & 4/4 \\
\midrule
\multicolumn{3}{l}{\emph{scarto medio dei best}} & \textbf{1,82\%} &
\multicolumn{2}{l}{\emph{scarto medio delle medie}} & \multicolumn{2}{r}{2,49\%} \\
\bottomrule
\end{tabular}
\end{table}

\paragraph{Lettura d'insieme.} Lo scarto medio dall'ottimo è dell'$\mathbf{1{,}82\%}$;
sette istanze su dieci restano sotto l'$1{,}5\%$ e tre sotto lo $0{,}5\%$
(A-n45-k7 si ferma a due unità di costo dall'ottimo). Le deviazioni standard sono
contenute ($1{,}9$--$9{,}2$, cioè meno dell'$1\%$ del costo): l'algoritmo è
\emph{stabile}, non fortunato. Il numero medio di generazioni necessarie a
raggiungere il risultato migliore è compreso fra $10{,}8$ e $19{,}2$ su circa
$24$ generazioni totali: i miglioramenti continuano cioè fino a ridosso
dell'esaurimento del budget, senza stagnazione.

Le due istanze più difficili sono E-n101-k14 ($3{,}94\%$) e P-n101-k4
($5{,}58\%$), entrambe con $100$ clienti; segue A-n80-k10 ($2{,}72\%$). La
correlazione con la dimensione è evidente ma non è l'intera storia, come mostra
la Sezione~\ref{sec:critica}.

\section{Analisi critica}
\label{sec:critica}

\subsection{Curve di convergenza}

\begin{figure}[H]
\centering
\includegraphics[width=\textwidth]{figures/fig_convergenza.pdf}
\caption{Curve di convergenza su quattro istanze rappresentative (una per
gruppo): costo della migliore soluzione trovata in funzione delle valutazioni
consumate. In azzurro le $5$ esecuzioni, in rosso la loro media, tratteggiato
l'ottimo noto.}
\label{fig:convergenza}
\end{figure}

Le curve (Figura~\ref{fig:convergenza}) hanno una forma comune: una discesa
molto ripida nelle prime $20\,000$--$30\,000$ valutazioni --- dove agiscono
l'inizializzazione greedy e le prime discese di ricerca locale --- seguita da un
lungo miglioramento a gradini, con salti progressivamente più rari ma non nulli.
Tre osservazioni critiche.

\textbf{(a) Il budget è ben dimensionato ma non sovrabbondante.} Il numero medio di
valutazioni necessarie a raggiungere il migliore risultato va da $155\,000$
(B-n78-k10) a $260\,000$ (P-n101-k4), su $350\,000$ disponibili; e confrontando lo
scarto a $50\,000$ valutazioni con quello finale (P-n101-k4: dal $7{,}84\%$ al
$6{,}37\%$; E-n101-k14: dal $5{,}47\%$ al $4{,}42\%$) si vede che l'ultima parte del
budget contribuisce ancora in modo apprezzabile. \emph{Non} siamo in saturazione:
un budget maggiore darebbe verosimilmente risultati migliori --- informazione
scomoda ma necessaria per interpretare correttamente i numeri.

\textbf{(b) La dispersione fra esecuzioni si riduce nel tempo.} Le cinque curve
partono separate (domina l'inizializzazione casuale) e convergono in una banda
stretta: l'algoritmo è attratto verso una regione di qualità simile
indipendentemente dal punto di partenza, coerentemente con le piccole deviazioni
standard della Tabella~\ref{tab:risultati}.

\textbf{(c) Nessuna stagnazione prematura.} Dopo l'introduzione della soppressione
dei duplicati non si osservano più i lunghi plateau terminali delle prime versioni:
ogni curva migliora anche nell'ultimo quarto del budget.

\subsection{Che cosa fanno realmente gli operatori}

Una critica legittima a molti lavori su algoritmi bio-ispirati è che gli
operatori vengono \emph{descritti} ma non \emph{misurati}: si dà per scontato che
funzionino. Abbiamo quindi strumentato l'algoritmo con contatori (che non
consumano numeri casuali e non alterano i risultati, verificati identici) per
misurare l'attività effettiva di ciascun meccanismo
(Tabella~\ref{tab:operatori}).

\begin{table}[H]
\centering
\small
\caption{Attività degli operatori, media su $5$ esecuzioni. \emph{Morti aging}:
anticorpi rimossi per superamento di $\tau_B$. \emph{Duplicati}: anticorpi
declassati dalla soppressione dei duplicati. \emph{LS efficace}: invocazioni della
ricerca locale che hanno prodotto un miglioramento, sul totale. \emph{Cloni
$>k$}: percentuale di cloni la cui decodifica non ammette $\le k$ percorsi.}
\label{tab:operatori}
\begin{tabular}{lrrrrr}
\toprule
Istanza & generazioni & morti aging & duplicati & LS efficace & cloni $>k$ \\
\midrule
A-n45-k7   & 24,0 & 16,6 & 27,4 & 70,8 / 71,2 & 10,2\% \\
A-n60-k9   & 24,0 & 24,8 & 32,2 & 69,8 / 70,4 & 40,6\% \\
A-n80-k10  & 24,0 & 46,6 & 18,0 & 69,0 / 70,2 & 58,8\% \\
B-n56-k7   & 24,0 & 48,8 & 46,6 & 69,2 / 70,0 & 0,2\% \\
B-n66-k9   & 24,0 & 6,4  & 11,6 & 70,0 / 70,0 & 83,6\% \\
B-n78-k10  & 24,0 & 30,8 & 25,8 & 69,2 / 70,0 & 46,2\% \\
E-n76-k8   & 24,0 & 14,4 & 17,2 & 70,0 / 70,0 & 34,8\% \\
E-n101-k14 & 24,0 & 26,2 & 17,8 & 69,8 / 70,0 & 46,2\% \\
P-n50-k10  & 27,2 & 28,0 & 36,4 & 79,6 / 80,6 & 93,3\% \\
P-n101-k4  & 24,0 & 30,4 & 33,2 & 69,8 / 70,0 & 0,0\% \\
\bottomrule
\end{tabular}
\end{table}

I dati confermano che nessun operatore è decorativo. L'\textbf{aging} è attivo
ovunque (da $6$ a $49$ rimozioni per esecuzione); il caso di B-n66-k9, dove le
morti sono poche, non è un'anomalia ma una conferma del design: lì gli anticorpi
migliorano di continuo, l'età viene azzerata, e quindi pochi raggiungono $\tau_B$
--- esattamente il comportamento voluto. La \textbf{soppressione dei duplicati}
interviene decine di volte per esecuzione. La \textbf{ricerca locale} produce un
miglioramento nel $\sim\!99\%$ delle invocazioni, il che giustifica il costo in
valutazioni che le abbiamo dedicato.

La colonna più eloquente è l'ultima. La percentuale di cloni non conformi al
vincolo di flotta varia da $0\%$ a $93\%$: sono le istanze a capacità satura
(P-n50-k10, B-n66-k9) quelle in cui il vincolo sul numero di veicoli è
\emph{davvero} il vincolo attivo, e in cui la selezione feasibility-first sta
lavorando a pieno regime. Su P-n101-k4 e B-n56-k7, all'estremo opposto, il
vincolo è di fatto inattivo. Questa è la misura del perché le scelte descritte
nella Sezione~\ref{sec:difficolta} erano necessarie e non sovra-ingegnerizzate.

\subsection{Da dove viene lo scarto residuo? Una scomposizione esatta}

L'analisi più utile che abbiamo condotto risponde a una domanda precisa: lo
scarto residuo dall'ottimo è dovuto al fatto che \emph{ordiniamo male} i clienti
dentro i percorsi, o al fatto che li \emph{raggruppiamo male} fra i veicoli?

Presa la migliore soluzione trovata, si
ri-ottimizza in modo \emph{ottimo} l'ordine di visita all'interno di ciascun
percorso, lasciando invariato l'insieme dei clienti assegnati a ogni veicolo
(algoritmo di Held--Karp, esatto, applicabile perché i percorsi contengono al più
$13$ clienti nella quasi totalità dei casi; su P-n101-k4, i cui percorsi arrivano
a $30$ clienti, si è usata una ricerca locale multi-avvio, calibrata verificando
che non riesce a migliorare i percorsi \emph{ottimi} dell'istanza). Si ottiene
allora la scomposizione additiva
\[
  \underbrace{\text{costo trovato} - \text{BKS}}_{\text{scarto totale}}
  \;=\;
  \underbrace{\text{costo trovato} - \text{costo riordinato}}_{\text{perdita di ordinamento}}
  \;+\;
  \underbrace{\text{costo riordinato} - \text{BKS}}_{\text{perdita di assegnamento}} .
\]

\begin{table}[H]
\centering
\small
\caption{Scomposizione dello scarto dall'ottimo. La perdita di ordinamento è
\textbf{nulla su tutte e dieci le istanze}: ogni percorso prodotto
dall'algoritmo è già ottimo come TSP.}
\label{tab:scomposizione}
\begin{tabular}{lrrrrr}
\toprule
Istanza & clienti/percorso & costo trovato & costo riordinato & perdita ordin. & perdita assegn. \\
\midrule
A-n45-k7   & 6,3  & 1148 & 1148 & \textbf{0} & 2 \\
A-n60-k9   & 6,6  & 1358 & 1358 & \textbf{0} & 4 \\
A-n80-k10  & 7,9  & 1811 & 1811 & \textbf{0} & 48 \\
B-n56-k7   & 7,9  & 710  & 710  & \textbf{0} & 3 \\
B-n66-k9   & 7,2  & 1332 & 1332 & \textbf{0} & 16 \\
B-n78-k10  & 7,7  & 1239 & 1239 & \textbf{0} & 18 \\
E-n76-k8   & 9,4  & 743  & 743  & \textbf{0} & 8 \\
E-n101-k14 & 7,1  & 1109 & 1109 & \textbf{0} & 42 \\
P-n50-k10  & 4,9  & 705  & 705  & \textbf{0} & 9 \\
P-n101-k4  & 25,0 & 719  & 719  & \textbf{0} & 38 \\
\midrule
\multicolumn{4}{l}{\emph{totale}} & \textbf{0 (0\%)} & \textbf{188 (100\%)} \\
\bottomrule
\end{tabular}
\end{table}

Il risultato (Tabella~\ref{tab:scomposizione}) è netto e, va detto,
controintuitivo: \textbf{la perdita di ordinamento è esattamente zero su tutte e
dieci le istanze}. Ogni singolo percorso prodotto dal nostro algoritmo è già
ottimo come TSP. Il $100\%$ dello scarto residuo ($188$ unità di costo in totale)
è \emph{perdita di assegnamento}: sbagliamo \emph{quali} clienti mettere insieme,
mai \emph{in che ordine} servirli.

Questa scoperta ha una spiegazione strutturale precisa, che riguarda proprio le
scelte progettuali descritte sopra. Il vicinato 2-opt intra-percorso è esplorato
a \emph{scansione completa} (i percorsi sono corti, quindi ce lo possiamo
permettere), mentre i vicinati inter-percorso sono ristretti ai $20$ vicini più
prossimi; inoltre lo Split può ritagliare solo segmenti \emph{contigui} del giant
tour, il che limita strutturalmente i raggruppamenti raggiungibili. La nostra
capacità di \emph{ordinare} è quindi molto più forte della nostra capacità di
\emph{riassegnare}. È un'informazione preziosa: indica esattamente dove
investire per migliorare l'algoritmo, e ci dice che ogni ulteriore sforzo sulla
componente TSP sarebbe completamente sprecato.

\paragraph{Un dettaglio istruttivo.} Contando gli auto-incroci geometrici dei
percorsi (una soluzione 2-opt-ottima, con distanze euclidee reali, non può averne)
se ne trova \emph{uno solo} in tutte e dieci le soluzioni, in un percorso di
B-n78-k10. Sembrerebbe un difetto della ricerca locale; in realtà quel percorso è
\emph{dimostrabilmente} ottimo come TSP (verificato con Held--Karp, costo $84$).
L'incrocio sopravvive perché le distanze TSPLIB sono \emph{arrotondate all'intero}:
la mossa 2-opt che lo eliminerebbe accorcia la lunghezza euclidea reale ma ha
$\Delta \ge 0$ sul costo intero. Non è un errore dell'algoritmo, è una proprietà
della metrica.

\subsection{Un esperimento naturale: E-n101-k14 contro P-n101-k4}

Analizzando i dati abbiamo notato che le istanze E-n101-k14 e P-n101-k4 hanno
coordinate e domande \textbf{identiche}: sono lo stesso insieme di $100$ clienti,
e differiscono unicamente per la capacità dei veicoli ($112$ contro $400$),
quindi per il numero di percorsi ($14$ contro $4$) e per la loro lunghezza
($\sim\!7$ contro $\sim\!25$ clienti). È un esperimento naturale controllato che
ci viene offerto gratuitamente: l'unica variabile che cambia è la
\emph{granularità} della partizione.

I risultati sono $3{,}94\%$ contro $5{,}58\%$: l'istanza a percorsi lunghi è più
difficile. Ma --- alla luce della scomposizione precedente --- \emph{non} perché
i percorsi lunghi siano più difficili da ordinare (li ordiniamo comunque in modo
ottimo). La ragione è che con $4$ percorsi da $25$ clienti lo spazio delle
partizioni ammissibili è enormemente più vasto e più ``piatto'': spostare un
cliente da un percorso all'altro cambia pochissimo il costo, i gradienti che
guidano la ricerca locale sono deboli, e la decodifica Split --- che può produrre
solo segmenti contigui del giant tour --- fatica a esplorare raggruppamenti
radicalmente diversi. La Figura~\ref{fig:evsp} lo mostra visivamente: i nostri quattro percorsi lunghi
(b) si intrecciano, mentre l'ottimo (c) li dispone in quattro regioni nettamente
separate --- eppure entrambe le soluzioni hanno percorsi TSP-ottimi.

\begin{figure}[H]
\centering
\includegraphics[width=0.96\textwidth]{figures/fig_panorama.pdf}
\caption{Un esperimento naturale. Le istanze E-n101-k14 \textbf{(a)} e P-n101-k4
\textbf{(b)} condividono \emph{esattamente} la stessa geometria (coordinate e
domande identiche) e differiscono solo per la capacità dei veicoli ($112$ contro
$400$): $14$ percorsi brevi contro $4$ percorsi lunghi. In \textbf{(c)} l'ottimo
di P-n101-k4: i singoli percorsi sono ottimi come TSP sia in (b) sia in (c); ciò
che cambia è la \emph{ripartizione} dei clienti, che nell'ottimo separa il piano
in regioni molto più nette.}
\label{fig:evsp}
\end{figure}

\subsection{Opinioni personali sui risultati}

Il risultato che considero più significativo non è lo scarto medio dell'$1{,}82\%$,
ma il fatto che le scelte progettuali più importanti siano emerse dai
\emph{fallimenti}, e non dalla teoria. La soppressione dei duplicati, la selezione
feasibility-first, il tetto di valutazioni per discesa: nessuna di queste tre era
prevista nella progettazione iniziale, e tutte e tre sono state imposte da
comportamenti anomali osservati sperimentalmente. In una metaeuristica, l'interazione fra gli
operatori genera comportamenti che la descrizione dei singoli operatori non lascia
prevedere, e solo la misurazione sistematica li rivela.

Sono soddisfatto della qualità raggiunta, consapevole che essa poggia in larga
parte sull'ibridazione: la ricerca locale produce la maggior parte della qualità,
mentre il motore immunologico fornisce la diversificazione che le impedisce di
restare intrappolata. Un algoritmo immunologico \emph{puro} non sarebbe
competitivo; ma neppure una ricerca locale multi-avvio, priva
dell'intensificazione attorno alle soluzioni élite, raggiungerebbe questi valori
(una singola discesa da Clarke--Wright si ferma fra il $2\%$ e l'$8{,}5\%$ di
scarto). Il limite più chiaro resta quello identificato dalla scomposizione: la
debolezza nella \emph{riassegnazione} dei clienti fra veicoli --- un limite
strutturale della rappresentazione a giant tour, il prezzo dell'eleganza della
decodifica Split, non un difetto di implementazione.

\section{Conclusioni}
\label{sec:conclusioni}

Il progetto ha prodotto un algoritmo immunologico a selezione clonale ibridato
con ricerca locale che raggiunge uno scarto medio dell'$1{,}82\%$ dagli ottimi
noti sulle dieci istanze richieste, con esecuzioni interamente riproducibili,
tutte le città servite e il vincolo sul numero di veicoli sempre rispettato.

Sul piano tecnico, ritengo che la scelta più felice sia stata quella della
rappresentazione a giant tour con decodifica esatta: essa ha reso banale
l'ammissibilità rispetto alla capacità, ha semplificato tutti gli operatori e ha
permesso di concentrare la ricerca su un unico spazio (le permutazioni). Il suo
prezzo --- la difficoltà nel riassegnare i clienti fra veicoli --- è emerso solo
alla fine, dalla scomposizione esatta dello scarto, ed è oggi l'indicazione più
precisa su dove intervenire.

Se dovessi proseguire il lavoro, agirei esattamente lì, e non altrove: vicinati
inter-percorso più potenti (catene di espulsione, scambi $\lambda$-interchange),
liste di candidati più ampie, o un operatore di mutazione che agisca
esplicitamente sui \emph{gruppi} anziché sull'ordine. Sarebbe inoltre naturale
esplorare la variante con \emph{penalty function} (opzione~5 della consegna): il
lavoro fatto sul vincolo di flotta mostra che le istanze a capacità satura
costringono la ricerca a muoversi lungo il bordo della regione ammissibile, ed è
plausibile che consentirne l'attraversamento controllato --- accettando
temporaneamente soluzioni con più di $k$ percorsi, penalizzate --- offrirebbe
scorciatoie che oggi ci sono precluse.

Sul piano metodologico, l'esperienza mi lascia una convinzione precisa: la
differenza fra un algoritmo che funziona e uno che non funziona non sta quasi mai
nella scelta del paradigma, ma nei dettagli di interazione fra gli operatori ---
dettagli che si scoprono solo misurando, e che nessuna descrizione teorica avrebbe
potuto suggerire in anticipo.

{\setlength{\parskip}{0pt}


\end{document}